\documentclass[10pt]{article}
\usepackage[T1]{fontenc}
\usepackage[utf8]{inputenc}
\usepackage[french]{babel}
\usepackage{amsmath,amssymb,mathrsfs,mathpazo}
\usepackage{geometry}
\geometry{a4paper,margin=2cm}
\usepackage{graphicx,array,tabularx,tabularray,colortbl,multirow,multicol,booktabs,enumitem,pifont,ulem,fancybox,diagbox,cancel}
\usepackage{pstricks,pstricks-add,pst-plot,pst-tree,pst-node,pst-text,pst-func,pst-math,pst-bspline,pst-3dplot,pst-eucl,pst-solides3d}
\let\pscalculate\undefined
\usepackage{pst-bezier}
\usepackage{tikz,pgfplots,tkz-tab}
\usetikzlibrary{arrows,arrows.meta,patterns,calc,positioning,trees,decorations.markings,shapes.geometric,shapes.misc,shapes.arrows,intersections,decorations.pathmorphing,decorations.pathreplacing,backgrounds,shadows,babel}
\pgfplotsset{compat=1.18}
\usepgfplotslibrary{fillbetween}
\usepackage[np]{numprint}
\usepackage{listings,makecell,pas-tableur,twoopt}
\let\barre\undefined
\usepackage[pstricks]{tablvar}
\definecolor{darkgreen}{RGB}{0,100,0}
\providecommand{\red}{\color{red}}
\providecommand{\blue}{\color{blue}}
\providecommand{\green}{\color{green}}
\providecommand{\black}{\color{black}}
\providecommand{\white}{\color{white}}
\providecommand{\gray}{\color{gray}}
\providecommand{\cyan}{\color{cyan}}
\providecommand{\magenta}{\color{magenta}}
\providecommand{\yellow}{\color{yellow}}
\providecommand{\violet}{\color{violet}}
\providecommand{\orange}{\color{orange}}
\providecommand{\euro}{\texteuro}
\providecommand{\EUR}{\texteuro}
\providecommand{\ieme}{\textsuperscript{e}}
\providecommand{\degres}{\ensuremath{^\circ}}
\providecommand{\degree}{\ensuremath{^\circ}}
\providecommand{\vv}[1]{\overrightarrow{#1}}
\providecommand{\Coord}[1]{#1}
\providecommand{\decosix}{}
\providecommand{\starredbullet}{\ensuremath{\bullet}}
\pagestyle{empty}
\setlength{\parindent}{0pt}
\usepackage{adjustbox}\DeclareUnicodeCharacter{2236}{:}\DeclareUnicodeCharacter{2190}{\ensuremath{\leftarrow}}
\providecommand{\R}{\mathbb{R}}
\providecommand{\N}{\mathbb{N}}
\providecommand{\Z}{\mathbb{Z}}
\providecommand{\Q}{\mathbb{Q}}
\providecommand{\D}{\mathbb{D}}
\providecommand{\C}{\mathbb{C}}
\providecommand{\e}{{\mathrm{e}}}
\providecommand{\vect}[1]{\overrightarrow{#1}}
\providecommand{\vectt}[1]{\overrightarrow{\mathrm{#1}}}
\providecommand{\vvt}[1]{\overrightarrow{\mathrm{#1}}}
\providecommand{\np}[2][]{{#2~#1}}
\providecommand{\vv}[1]{{\overrightarrow{#1}}}
\renewcommand{\d}{{\mathrm{d}}}
\providecommand{\degres}{{^\circ}}
\providecommand{\degre}{{^\circ}}
\providecommand{\blue}{\color{blue}}
\providecommand{\red}{\color{red}}
\providecommand{\enskip}{\quad}
\providecommand{\llbracket}{\lbrack\!\lbrack}
\providecommand{\rrbracket}{\rbrack\!\rbrack}
\providecommand{\euro}{{\text{€}}}
\providecommand{\no}{{\text{n°}}}
\providecommand{\interval}[2]{{[#1;#2]}}
\providecommand{\intervFO}[2]{{[#1;#2[}}
\providecommand{\dotfill}{\ldots}

\providecommand{\ds}{\displaystyle}
\providecommand{\textding}[1]{#1}
\providecommand{\Oij}{\(\mathrm{O};\vec{\imath},\vec{\jmath}\)}
\providecommand{\Oijk}{\(\mathrm{O};\vec{\imath},\vec{\jmath},\vec{k}\)}
\providecommand{\Ouv}{\(\mathrm{O};\vec{u},\vec{v}\)}
\providecommand{\intervFF}[2]{\ensuremath{[#1;#2]}}
\renewcommand{\euro}{\ifmmode\text{\texteuro}\else\texteuro\fi}
\DeclareUnicodeCharacter{00B0}{\ifmmode^{\circ}\else\textdegree\fi}
\DeclareUnicodeCharacter{20AC}{\ifmmode\text{\texteuro}\else\texteuro\fi}
% Énoncés : Éducation nationale. Transcription et corrigés : APMEP, Denis Vergès et contributeurs. Découpage et classement : CoopMaths, Sébastien Lozano, Éric Elter et contributeurs. Adaptation de la présentation : Maths Explorer.
% CC BY-SA : https://coopmaths.fr/www/about/
% https://coopmaths.fr/alea/static/bac/2023/tex/bac_2023_09_sujet1_ameriquesud_4.tex
\begin{document}
\noindent\textit{APMEP / CoopMaths. CC BY-SA. Présentation adaptée par Maths Explorer.}\par\medskip
\medskip

\textbf{Partie A}

\medskip

Le but de la partie A est d'étudier le comportement de la suite $\left(u_n\right)$ définie par $u_0 =  0,3$ et par la relation de récurrence, pour tout entier naturel $n$ :
\[u_{n+1} =  2u_n\left(1 - u_n\right).\]

Cette relation de récurrence s'écrit $u_{n+1} =  f\left(u_n\right)$, où $f$ est la fonction définie sur $\R$ par : 
\[f(x) = 2x\left (1 - x\right ).\]

\smallskip
\begin{enumerate}
\item Démontrer que la fonction $f$ est strictement croissante sur l'intervalle $\left[0~;~\dfrac12\right]$.
\item On admet que pour tout entier naturel $n$, $0 \leqslant u_n \leqslant \dfrac12$.

Calculer $u_1$ puis effectuer un raisonnement par récurrence pour démontrer que pour tout rentier naturel $n,\: u_n \leqslant u_{n+1}$.
\item En déduire que la suite $\left(u_n\right)$ est convergente.
\item Justifier que la limite de la suite $\left(u_n\right)$ est égale à $\dfrac12$.
\end{enumerate}

\bigskip

\textbf{Partie B}

\medskip

Le but de cette partie est d'étudier un modèle d'évolution d'une population.

En 2022, cette population compte \np{3000} individus.

\smallskip

On note $P_n$ l'effectif en milliers de la population l'année $2022 + n$. Ainsi $P_0 = 3$.

Selon un modèle inspiré du modèle de Verhulst, mathématicien belge du XIX\up{e} siècle, on considère que, pour tout entier naturel $n$ :

\begin{center}$P_{n+1} -  P_n = P_n\left(1 -  b \times P_n\right)$, où $b$ est un réel strictement positif.\end{center}

Le réel $b$ est un facteur de freinage qui permet de tenir compte du caractère limité des
ressources du milieu dans lequel évoluent ces individus.

\medskip

\begin{enumerate}
\item Dans cette question $b  = 0$.
	\begin{enumerate}
		\item Justifier que la suite $\left(P_n\right)$ est une suite géométrique dont on précisera la raison.
		\item Déterminer la limite de $P_n$.
	\end{enumerate}
\item Dans cette question $b = 0,2$.
	\begin{enumerate}
		\item Pour tout entier naturel $n$, on pose $v_n = 0,1 \times P_n$.

Calculer $v_0$ puis montrer que, pour tout entier naturel $n$,\:$v_{n+1} = 2v_n\left(1 - v_n \right)$.
		\item Dans ce modèle, justifier que la population se stabilisera autour d'une valeur que l'on précisera.
	\end{enumerate}
\end{enumerate}

\end{document}
